\scanpage{28}
\begin{center}4 - 22 -\end{center}
\hfill n° 392

PROPOSITION 10.- (i) \uline{Soit} $0\longrightarrow T_1\longrightarrow T_2\longrightarrow T_3\longrightarrow0$ \uline{une suite exacte de} $A$-\uline{modules de torsion de type fini}. \uline{On a alors} :
\[\chi(T_2)=\chi(T_1)+\chi(T_3).\]
(ii) \uline{S'il existe un pseudo-isomorphisme} $f:T_1\longrightarrow T_2$, \uline{on a} $\chi(T_1)=\chi(T_2)$.

(iii) \uline{Pour que} $\chi(T)=0$, \uline{il faut et il suffit que} $T$ \uline{soit pseudo-nul}.

Vu la déf. 4, il suffit de considérer pour chaque $\uline{p}\in P$, les valeurs de $\ell_{\uline{p}}$ pour les modules de torsion considérés. La propriété (i) résulte alors du chap. II, § 2, n° 4, th. 1 et de l'additivité des longueurs dans une suite exacte, (\uline{Alg.}, chap. II, 3e éd., § 1, n° 10, prop. 16), et les propriétés (ii) et (iii) résultent aussitôt des définitions du n° 4.

COROLLAIRE.- \uline{Soit} $0\longrightarrow T_n\longrightarrow T_{n-1}\longrightarrow\cdots\longrightarrow T_0\longrightarrow0$ \uline{une suite exacte de} $A$-\uline{modules de torsion de type fini}. \uline{On a} $\sum_{i=0}^n(-1)^i\chi(T_i)=0$.

Vu le chap. II, § 2, n° 4, th. 1, cela résulte encore de la propriété analogue des $\ell_{\uline{p}}$, (\uline{Alg.}, chap. II, 3e éd., § 1, n° 10, cor. 3 de la prop. 16).

Rappelons, (chap. II, § 5, n° 4), que l'on peut parler de \uline{l'ensemble} $\uline{F}(A)$ \uline{des classes de} $A$-\uline{modules de type fini} pour la relation d'isomorphie pour tout $A$-module $M$ de type fini, on désigne par $c\ell(M)$ l'élément correspondant de $\uline{F}(A)$ ; nous désignerons par $\uline{T}(A)$ la partie de $\uline{F}(A)$ formée des classes de $A$-modules de torsion de type fini. Il est clair que $\chi$ définit une application de $\uline{T}(A)$ dans $D(A)$, notée encore $\chi$, telle que $\chi(c\ell(T))=\chi(T)$.

PROPOSITION 11.- \uline{Soient} $G$ \uline{un groupe commutatif}, \uline{noté additivement}, \uline{et} $\varphi:\uline{T}(A)\longrightarrow G$ \uline{une application} ; \uline{pour tout} $A$-\uline{module de torsion de type fini} $T$, \uline{on pose encore}, \uline{par abus de langage}, $\varphi(T)=\varphi(c\ell(T))$. \uline{On suppose vérifiées les conditions suivantes} :

1) Si $0\longrightarrow T_1\longrightarrow T_2\longrightarrow T_3\longrightarrow0$ \uline{est une suite exacte de} $A$-\uline{modules de torsion de type fini}, \uline{on a} $\varphi(T_2)=\varphi(T_1)+\varphi(T_3)$.
\mspage{27}
